\section{Rectangular tests from subspace expansion}
\label{sec:expansion}

Fix the variable and formula-size bounds $n,s\geq1$. We construct
rectangular matrices $T_0,\ldots,T_n$ such that
$\sum_{i=0}^n A_i\otimes T_i$ is injective whenever
$A_0+\sum_{i=1}^n A_i x_i$ is a feasible affine pencil of order
at most $2s+1$. The construction uses an indeterminate $t$; no
specialization of a formal Laurent series will be needed. Set
\begin{equation}\label{eq:parameters}
\begin{gathered}
w=n+1,\qquad B=2s+1,\qquad \Delta=2w^3B,\\
M=\min\{2^a:a\in\mathbb Z_{\geq0},\ 2^a>B\Delta\},
\qquad N=M-\Delta.
\end{gathered}
\end{equation}
The bound $B$ will accommodate all certificate pencils in
Proposition~\ref{prop:formula-certificates}. The padding width $\Delta$
provides the expansion margin, while $M>B\Delta$ makes the eventual
rank deficit smaller than one block. These choices give
$1\leq N<M\leq2B\Delta$.

For a field $K$ and an integer $d\geq 1$, let
\[
  \mathcal P_d(K)=\{f\in K[z]:\deg f<d\},
\]
with its monomial basis.  Put $F=\Q((t))$ and $\mathcal O=\Q[[t]]$.
For $0\leq i\leq n$, define an $F$-linear map
$T_i\colon\mathcal P_N(F)\longrightarrow\mathcal P_M(F)$ by
\begin{equation}
  f=\sum_{a=0}^{N-1}c_{ia}(f)(z-i)^a,
  \qquad
  T_i f=\sum_{a=0}^{N-1}c_{ia}(f)t^a(z-i)^{M-1-a}.
  \label{eq:rectangular-test}
\end{equation}
The exponents in the output are nonnegative and distinct.  Since each
$t^a$ is nonzero in $F$, every $T_i$ is injective.  Its matrix in the
monomial bases has entries in $\mathbb Z[t]$.  More explicitly, its column
corresponding to $z^b$, for $0\leq b<N$, is the coefficient vector of
\begin{equation}
  T_i(z^b)=\sum_{a=0}^{b}\binom{b}{a}i^{b-a}t^a(z-i)^{M-1-a}.
  \label{eq:test-column}
\end{equation}
Thus all its entries have $t$-degree at most $N-1$.  Let $R_i(t)$ be the
$M\times M$ matrix obtained by appending $M-N$ zero columns to $T_i$.

The conclusion of this section is the following proposition.

\begin{proposition}[Rectangular rank bound]
\label{prop:rectangular}
Let $1\leq q\leq B$ and $A_0,\ldots,A_n\in\Mat_q(\Q)$.  Suppose that,
for some extension field $E/\Q$, some integer $e\geq1$, and some
$X_1,\ldots,X_n\in\Mat_e(E)$, the matrix
\[
  A_0\otimes I_e+\sum_{i=1}^n A_i\otimes X_i
\]
is invertible.  Then $\sum_{i=0}^n A_i\otimes T_i$ is injective over $F$,
and
\begin{equation}
  \rank_{\Q(t)}\left(\sum_{i=0}^n A_i\otimes R_i(t)\right)
  \geq qN>(q-1)M.
  \label{eq:padded-rank}
\end{equation}
\end{proposition}

We prove the proposition by showing that the maps $T_i$ jointly expand
every nonzero subspace of dimension at most $B$.  Two elementary facts will
control the vanishing orders of their leading coefficients.
The ordinary-Wronskian multiplicity method used in the subspace designs
of Guruswami and Kopparty~\citep[Section~5.1 of the preprint]{GK16}
is a predecessor of this argument. Their construction bounds
intersections with multiplicity subspaces; here Taylor reversal
relates the vanishing sequences of the input and image spaces.
Forbes and Guruswami~\citep{FG15} explain broader connections among
subspace designs, rank condensers, and dimension expansion.
The lossless and unbalanced dimension expanders of Guruswami, Resch,
and Xing~\citep{GRX18} provided a comparison for the required
small-subspace expansion parameters during development.  Their
linearized-polynomial construction is not used here.
We prove the order estimate and the rectangular expansion statement
directly for the maps above.

\subsection{Vanishing orders and integral reduction}

Let $V\subseteq\mathcal P_d(\Q)$ have dimension $h\geq1$.  Row reduction
of Taylor coefficient vectors about $z=i$ gives a basis with distinct
vanishing orders
\[
  0\leq d_{i1}<\cdots<d_{ih}<d.
\]
These orders describe the entire vanishing filtration: for every integer
$r\geq0$,
\begin{equation}
  \dim_{\Q}\{f\in V:(z-i)^r\text{ divides }f\}
  =\#\{j:d_{ij}\geq r\}.
  \label{eq:vanishing-filtration}
\end{equation}
Indeed, in a nonzero linear combination of the echelon basis, the first
basis vector with a nonzero coefficient supplies its lowest Taylor term.

The leading-term Wronskian computation in the following lemma is
classical; the echelon-basis and Vandermonde argument is also presented
by Bostan and Dumas~\citep[Lemmas~1--3]{BD10}.

\begin{lemma}[Wronskian order bound]
\label{lem:wronskian-orders}
With the notation above, at the $w=n+1$ points $0,\ldots,n$ one has
\begin{equation}
  \sum_{i=0}^n\sum_{j=1}^h d_{ij}
  \leq h(d-1)+w\binom{h}{2}.
  \label{eq:wronskian-bound}
\end{equation}
\end{lemma}

\begin{proof}
For a basis $f_1,\ldots,f_h$ of $V$, consider its Wronskian
\[
  \mathcal W(z)=\det\bigl(f_j^{(r)}(z)\bigr)_{
       0\leq r<h,\,1\leq j\leq h}.
\]
Changing the basis multiplies $\mathcal W$ by a nonzero rational
constant.  At a fixed point $i$, use the echelon basis and write
$f_j=\alpha_j(z-i)^{d_{ij}}+$ higher powers, where $\alpha_j\neq0$.
The contribution of these leading terms to the determinant is
\[
  \left(\prod_{j=1}^h\alpha_j\right)
  \det\bigl((d_{ij})_r\bigr)_{
       0\leq r<h,\,1\leq j\leq h}
  (z-i)^{\sum_jd_{ij}-\binom h2},
\]
where $(X)_r=X(X-1)\cdots(X-r+1)$ and $(X)_0=1$.  Because these are
monic polynomials of successive degrees, the displayed coefficient
determinant is the Vandermonde
$\prod_{a<b}(d_{ib}-d_{ia})$, which is nonzero.  Terms using a higher
power from any $f_j$ have strictly larger total exponent, or vanish on
differentiation.  The displayed exponent is nonnegative because the
$d_{ij}$ are distinct nonnegative integers.  Consequently
\[
  \ord_{z=i}\mathcal W
  =\sum_{j=1}^h d_{ij}-\binom h2,
\]
and in particular $\mathcal W$ is not the zero polynomial.

Each entry in its determinant has degree at most $d-1$, so
$\deg\mathcal W\leq h(d-1)$.  The sum of its vanishing orders at the
distinct points $0,\ldots,n$ is at most its degree.  This gives
\eqref{eq:wronskian-bound}.
\end{proof}

For a subspace $U\subseteq F^D$, define its reduction by
\[
  \overline U
  =\{u\bmod t:u\in U\cap\mathcal O^D\}\subseteq\Q^D.
\]
It is important here to use all integral vectors of the subspace, not
merely the reductions of an arbitrarily chosen basis.

\begin{lemma}[Integral bases of Laurent-series subspaces]
\label{lem:integral-reduction}
Every $k$-dimensional subspace $U\subseteq F^D$ has a basis
$u_1,\ldots,u_k\in\mathcal O^D$ whose reductions are linearly independent.
For such a basis,
\begin{equation}
  U\cap\mathcal O^D=\bigoplus_{j=1}^k\mathcal O u_j,
  \qquad
  \dim_{\Q}\overline U=k.
  \label{eq:integral-lattice}
\end{equation}
\end{lemma}

\begin{proof}
We first construct the basis by induction on $k$, the case $k=0$ being
immediate.  Multiply a nonzero vector of $U$ by a power of $t$ so that
all its coordinates are integral and some coordinate is a unit.  Call
this vector $u_1$ and its unit coordinate $c$.  The subspace of $U$
whose $c$-th coordinate is zero has dimension $k-1$.  By induction it
has an integral basis with independent reductions.  All those
reductions have zero $c$-th coordinate, whereas the reduction of $u_1$
does not.  Together they are therefore independent and form a basis
of $U$ over $F$.

Some $k\times k$ coordinate minor of the resulting basis matrix is a
unit in $\mathcal O$.  If an $F$-linear combination of its columns is
integral, restriction to those coordinates and inversion of that minor
show that all its coefficients belong to $\mathcal O$.  This proves
the first assertion of \eqref{eq:integral-lattice}; reducing its
coefficients modulo $t$ proves the second.
\end{proof}

For a nonzero vector $f$ with Laurent-series coordinates, let
$\nu_t(f)$ denote the minimum of their $t$-valuations.  Its
\emph{leading reduction} is the nonzero vector
$t^{-\nu_t(f)}f\bmod t$.  Any invertible coordinate change over $\Q$
preserves both the integral lattice and this minimum valuation: the
change-of-basis matrix and its inverse both have entries in $\mathcal O$.
In particular, this applies to conversion between monomial and Taylor
coordinates.  All these statements concern formal series and require
no convergence assumption.

\subsection{Expansion of every small subspace}

For an integral input having a unit Taylor coefficient with index $a\leq A$,
the leading reduction of its image uses only indices $a\leq A$; reversal
turns these into output orders at least $M-1-A$. The Wronskian bound then
rules out too small a common image space. The proof below applies this
observation to a suitable subspace for each vanishing order.

\begin{lemma}[Small-subspace expansion]
\label{lem:small-expansion}
For every $k$-dimensional subspace $U\subseteq\mathcal P_N(F)$ with
$1\leq k\leq B$,
\begin{equation}
  \dim_F\left(\sum_{i=0}^n T_iU\right)>(w-1)k.
  \label{eq:small-expansion}
\end{equation}
\end{lemma}

\begin{proof}
Put $W=\sum_{i=0}^nT_iU$ and $\ell=\dim_FW$.  Since $T_0$ is
injective, $\ell\geq k$.  By Lemma~\ref{lem:integral-reduction},
$\overline U\subseteq\mathcal P_N(\Q)$ has dimension $k$ and
$\overline W\subseteq\mathcal P_M(\Q)$ has dimension $\ell$.
Write their increasing vanishing orders at $i$ as
\[
  a_{i1}<\cdots<a_{ik},
  \qquad
  b_{i1}<\cdots<b_{i\ell}.
\]
We claim that, for every $i$ and every $1\leq j\leq k$,
\begin{equation}
  b_{i,\ell-j+1}\geq M-1-a_{ij}.
  \label{eq:order-reversal}
\end{equation}

Fix $i,j$ and set $A=a_{ij}$.  Projection of $\overline U$ onto its
Taylor coefficients of orders $0,\ldots,A$ has rank $j$.  Choose $j$
vectors in $\overline U$ whose projections are independent, lift them
to integral vectors of $U$, and let $U'$ be their $F$-span.  Their
reductions are independent, so Lemma~\ref{lem:integral-reduction}
shows that these reductions form a basis of $\overline{U'}$.  The chosen Taylor
projection is therefore injective on $\overline{U'}$.

Take a nonzero $f\in U'$ and multiply it by a power of $t$ so that
$\nu_t(f)=0$.  All its Taylor coefficients $c_{ia}(f)$ are integral.
Because its nonzero reduction has a nonzero low Taylor projection,
at least one coefficient with $a\leq A$ is a unit.  With the convention
$\nu_t(0)=+\infty$, put
\[
  m=\min_{0\leq a<N}\bigl(\nu_t(c_{ia}(f))+a\bigr).
\]
Then $m\leq A$, whereas for every $a>A$ one has
$\nu_t(c_{ia}(f))+a\geq a>A\geq m$.  Thus the leading reduction of
$T_i f$ uses only output powers $(z-i)^{M-1-a}$ with $a\leq A$.
These powers are independent, so their leading coefficients cannot
all cancel.  Since Taylor coordinate changes preserve minimum
valuation, $m$ is also the minimum valuation of the output in monomial
coordinates.  Its leading reduction consequently vanishes at $i$
to order at least $M-1-A$. Figure~\ref{fig:taylor-reversal} illustrates this reversal.

\begin{figure}[t]
\centering
\begin{tikzpicture}[>=Latex,font=\small,x=1cm,y=1cm]
\fill[blue!8] (0,-.23) rectangle (3.3,.23);
\fill[blue!8] (3.0,-2.23) rectangle (6.3,-1.77);
\draw[->] (-.2,0)--(6.6,0) node[right] {$a$};
\draw[->] (-.2,-2)--(6.6,-2) node[right] {order in $z-i$};
\foreach \x/\lab in {0/0,3.3/A,6.3/{N-1}} {
 \draw (\x,-.08)--(\x,.08);
 \node[above=4pt] at (\x,0) {$\lab$};
}
\draw (0,-2.08)--(0,-1.92);
\node[below=4pt] at (0,-2) {$\Delta$};
\draw (3,-2.08)--(3,-1.92);
\node[below=4pt] at (3,-2) {$M-1-A$};
\draw (6.3,-2.08)--(6.3,-1.92);
\node[below=4pt] at (6.3,-2) {$M-1$};
\draw[->,blue!60!black,thick] (.5,-.3)--(5.8,-1.68);
\draw[->,blue!60!black,thick] (2.8,-.3)--(3.5,-1.68);
\node[anchor=south,font=\footnotesize] at (1.65,.45) {possible surviving indices};
\node[align=center,anchor=south] at (3.15,.95) {A unit coefficient with $a\leq A$ forces minimum $t$-valuation $\leq A$.};
\node[anchor=south] at (3.2,1.5) {$c_a(z-i)^a\ \longmapsto\ c_at^a(z-i)^{M-1-a}$};
\end{tikzpicture}
\caption{Taylor reversal in the expansion argument (schematic, not to scale). For a normalized input with integral Taylor coefficients and a unit among $a\leq A$, every term surviving in the leading reduction of its image has $a\leq A$. The reduction therefore vanishes to order at least $M-1-A$ at $i$. The proof applies this to every vector in a $j$-dimensional subspace, with $A=a_{ij}$.}
\label{fig:taylor-reversal}
\end{figure}


This conclusion controls the whole reduction $\overline{T_iU'}$.
Indeed, every nonzero $v\in T_iU'$ has a preimage $f\in U'$.  Normalizing
$f$ multiplies $v$ by the same power of $t$, so the leading reduction
of $v$ is exactly the leading reduction already considered.  An integral
$v$ reduces either to this leading reduction or to zero.  We have not
assumed that $T_i(U'\cap\mathcal O^N)$ is the full integral lattice of
$T_iU'$.  By injectivity of $T_i$ and
Lemma~\ref{lem:integral-reduction}, $\overline{T_iU'}$ has dimension $j$.
It lies in $\overline W$ and in its vanishing filtration of order at
least $M-1-A$.  Formula~\eqref{eq:vanishing-filtration} now gives
\eqref{eq:order-reversal}.

For each fixed $i$, summing \eqref{eq:order-reversal} over $j=1,\ldots,k$
uses the largest $k$ distinct orders of the same space $\overline W$.
The remaining orders are nonnegative.  Although the auxiliary spaces
$U'$ may depend on $j$, these are bounds on the fixed filtration of
$\overline W$ and can therefore be summed.  Applying
Lemma~\ref{lem:wronskian-orders} to both reduced spaces gives
\begin{align*}
  \ell(M-1)+w\binom\ell2
  &\geq\sum_{i=0}^n\sum_{h=1}^{\ell}b_{ih}\\
  &\geq wk(M-1)-\sum_{i=0}^n\sum_{j=1}^k a_{ij}\\
  &\geq wk(M-1)-k(N-1)-w\binom k2\\
  &=(w-1)k(M-1)+k\Delta-w\binom k2.
\end{align*}
If $\ell\leq(w-1)k$, this implies
\[
  k\Delta\leq w\left(\binom k2+\binom\ell2\right)
  \leq\frac w2\bigl(1+(w-1)^2\bigr)k^2
  \leq\frac{w^3}{2}Bk
  <2w^3Bk=k\Delta,
\]
a contradiction.  This proves the strict expansion bound.
\end{proof}

\subsection{From expansion to the rank bound}

The proof uses the standard shrunk-subspace obstruction: if every
coefficient matrix sends a common subspace into one of smaller
dimension, no matrix substitution can be invertible.
This is the elementary obstruction direction of the
noncommutative-rank characterization of Fortin and
Reutenauer~\citep[Section~2, Theorem~1]{FR04}; see also
\citet[Definition~1.1 and Section~5.1]{IQS18}.
We derive the obstruction from expansion and give the required
extension-field comparison explicitly.

\begin{proof}[Proof of Proposition~\ref{prop:rectangular}]
Suppose that $\sum_iA_i\otimes T_i$ has a nonzero kernel vector,
written as $q$ components $f_1,\ldots,f_q\in\mathcal P_N(F)$.  Let
$U$ be their span, of dimension $1\leq k\leq q\leq B$, and choose
a basis $u_1,\ldots,u_k$ of $U$.  Write
\[
  f_\alpha=\sum_{j=1}^k P_{\alpha j}u_j,
  \qquad P\in\Mat_{q\times k}(F).
\]
Since the $f_\alpha$ span $U$, the matrix $P$ has rank $k$.  Each
row of the concatenated matrix
\[
  C=[\,A_0P\ \ A_1P\ \ \cdots\ \ A_nP\,]
\]
is a relation among the $wk$ vectors $T_i u_j$: this is exactly the
corresponding component of the kernel equation.  Those vectors span
$\sum_iT_iU$, whose dimension is greater than $(w-1)k$ by
Lemma~\ref{lem:small-expansion}.  Their relation space therefore has
dimension less than $k$, and hence $\rank C<k$.

Let $V$ be the column space of $P$ and $V_-$ the column space of $C$,
both in $F^q$.  Equality of row rank and column rank yields
\begin{equation}
  \dim_F V=k,\qquad \dim_FV_-<k,
  \qquad A_iV\subseteq V_-\quad(0\leq i\leq n).
  \label{eq:shrunk-subspace}
\end{equation}
Thus every coefficient matrix sends the same input subspace into a
strictly smaller common output subspace.

To compare this obstruction with the given witness over $E$, introduce
a fresh formal variable $\tau$ and work in $J=E((\tau))$.  Embed $E$
as the constant series and embed $F$ by
\[
  \sum_m q_m t^m\longmapsto\sum_m q_m\tau^m.
\]
These embeddings agree on $\Q$, which contains all the $A_i$.  A
transcendental element already in $E$, even one named $t$, remains a
coefficient and is not identified with $\tau$.  The embeddings preserve
matrix ranks, since nonzero minors stay nonzero.  Denote the extended
subspaces in \eqref{eq:shrunk-subspace} by $V_J$ and $(V_-)_J$.
Putting $X_0=I_e$, the witness evaluation maps
\[
  \left(\sum_{i=0}^n A_i\otimes X_i\right)
       (V_J\otimes_J J^e)
  \subseteq (V_-)_J\otimes_J J^e.
\]
The domain on the left has dimension $ke$, and the space on the right
has dimension strictly less than $ke$.  This contradicts invertibility
of the witness matrix, whose inverse over $E$ remains an inverse over
$J$.  The rectangular evaluation is therefore injective.

In the padded evaluation $\sum_iA_i\otimes R_i(t)$, selecting the
first $N$ columns within each of the $q$ blocks recovers the rectangular
evaluation.  Its rank over $F$ is at least $qN$.  The corresponding
nonzero maximal minor belongs to $\Q[t]$, so it is also nonzero over
$\Q(t)$.  In particular, no coefficient of the formal subspaces or
their bases has to be specialized.  Finally,
\[
  qN-(q-1)M=M-q\Delta\geq M-B\Delta>0,
\]
which proves \eqref{eq:padded-rank}.
\end{proof}
